\chapter{Menghidupkan Kembali Warisan Bahasa di Era Digital}

\section{Bahasa Sangihe: Mutiara Kebudayaan di Tapal Batas Utara}
Di ujung utara kepulauan Nusantara, terbentang gugusan pulau yang indah di wilayah Nusa Utara (Kabupaten Kepulauan Sangihe, Sulawesi Utara). Di tanah berombak inilah Bahasa Sangihe (\textit{Bahasa Sangir}) lahir, tumbuh, dan dituturkan secara turun-temurun oleh para pelaut ulung, penutur sastra lisan, dan masyarakat adat. Bahasa ini sarat dengan kearifan lokal, filosofi kehidupan bahari, serta karya sastra bernilai tinggi seperti wiracarita \textit{Sasahara} dan puisi epik \textit{Biong i Himbawo}.

Namun, seiring derasnya arus modernisasi dan peralihan generasi penutur muda ke bahasa pergaulan umum (Melayu Manado dan Bahasa Indonesia gaul), keberadaan bahasa daerah ini menghadapi tantangan yang nyata. Jika tidak ada upaya serius untuk mendokumentasikannya ke dalam bentuk digital dan menghadirkannya di gawai modern, bahasa yang kaya ini berisiko perlahan-lahan memudar dari percakapan sehari-hari.

\section{Tantangan AI: Mengapa Bahasa Daerah Kerap Tertinggal?}
Dalam beberapa tahun terakhir, dunia menyaksikan lompatan raksasa dalam teknologi kecerdasan buatan (\textit{Artificial Intelligence} / AI). Model-model bahasa besar (\textit{Large Language Model} atau LLM) seperti ChatGPT atau Qwen mampu menjawab pertanyaan rumit, menulis esai, bahkan menerjemahkan berbagai bahasa dunia dengan fasih.

\begin{infobox}[Analogi Perpustakaan Raksasa Dunia]
Bayangkan sebuah model AI modern seperti seorang sarjana pintar yang telah membaca miliaran buku dari perpustakaan terbesar dunia. Masalahnya, 95\% buku di perpustakaan tersebut ditulis dalam Bahasa Inggris, Mandarin, Spanyol, atau Bahasa Indonesia baku. Bahasa-bahasa daerah seperti Bahasa Sangihe hampir tidak pernah ia temukan di rak perpustakaan itu. Akibatnya, saat kita mengajaknya berbicara dalam Bahasa Sangihe, sang sarjana akan bingung atau sekadar menebak-nebak kata.
\end{infobox}

Dalam bidang linguistik komputasi, bahasa seperti Sangihe dikategorikan sebagai \textit{low-resource language} (bahasa berlindung sumber daya komputasi minim). Tantangan utamanya meliputi:
\begin{enumerate}
    \item \textbf{Kelangkaan Data Digital:} Tidak banyak situs web atau artikel daring yang memuat teks bahasa Sangihe dalam jumlah besar.
    \item \textbf{Kerumitan Tanda Diakritik:} Huruf vokal dalam bahasa Sangihe memiliki penanda khusus yang menentukan arti, yang kerap hilang saat diketik di ponsel pintar.
    \item \textbf{Morfologi yang Kaya:} Satu kata dasar bisa berubah wujud menjadi puluhan bentuk baru melalui awalan, sisipan di tengah kata, dan akhiran.
\end{enumerate}

\section{Misi Proyek: Membangun Mesin Penerjemah Cerdas}
Proyek ini lahir dengan sebuah cita-cita mulia namun terukur: \textbf{Membangun mesin penerjemah dua arah antara Bahasa Sangihe dan Bahasa Indonesia yang akurat, menghormati kaidah bahasa formal, dan dapat dijalankan di komputer biasa secara luring (offline) tanpa bergantung pada peladen cloud berbayar.}

Untuk mewujudkannya, proyek ini memadukan dua pendekatan teknologi mutakhir:
\begin{itemize}
    \item \textbf{Teknologi Hybrid RAG (\textit{Retrieval-Augmented Generation}):} Membekali sang model AI dengan "buku pintar" berupa kamus digital leksikon klasik dan bank ingatan kalimat paralel yang dapat ia buka setiap kali diminta menerjemahkan.
    \item \textbf{Pelatihan Khusus (\textit{Fine-Tuning}):} Melatih kembali saraf pemikiran model fondasi \texttt{\StatBaseModel} dengan puluhan ribu contoh kalimat terjemahan nyata agar ia memahami gaya, tata krama, dan jiwa Bahasa Sangihe.
\end{itemize}

Buku ini akan mengajak Anda menyusuri perjalanan lengkap bagaimana teknologi ini dirancang dari nol, langkah demi langkah, dengan penjelasan yang bersahabat bagi siapa saja.
